\section{Evidence from Baseline}\label{sec:empirics}

This section sets the incentive results beside transaction-level data from a deployed design. It is descriptive: the data come from one protocol over seven months, the quantities that the theory conditions on (horizons, outside rates) are not observed, and nothing here identifies a causal effect. Its main finding is a limit of the model. By amount, \ELeverAmtPct\% of the new debt in the sample was borrowed to buy more of the token, a use that \cref{sec:incentives,sec:equilibrium} do not contain.

\subsection{Setting and data}

Baseline's current system went live on 20 April 2026. Its documentation states the terms of its credit: a holder borrows the full floor value (``BLV'') of the locked tokens in the pool's reserve asset, pays a one-time origination fee, and faces no interest, no expiry and no liquidation \citep{baselinedocs,baselinefaq}; borrowers are staked automatically and are eligible for the same rewards as other stakers \citep{baselinestake}; and the sell fee is charged on the part of the price above the floor \citep{baselinefees}. These are the documented terms; we have not verified the deployed contracts, which are upgradeable. The log records \EUpgrades{} component upgrades on \EUpDates{} dates in the window: \EUpCredit{} of the credit component, \EUpSwap{} of the swap component, \EUpController{} of the controller, \EUpFactory{} of the factory and \EUpLens{} of the read-only lens. We do not know what each upgrade changed. Two checks bound the risk to the reconstruction. First, every event type that we decode has the same number of indexed fields and the same data length throughout the sample, except the carry-over event at launch, whose data are arrays. This check establishes stable topic counts and data lengths; it does not establish that the encoding or the meaning of a field remained unchanged. Second, the reconstructed totals equal the contract's own at the last block. The last upgrade of the credit component is at block \EUpLastCredit, and we report the main estimate on both sides of it. In the notation of \cref{sec:general} the documented terms correspond to $\gbar=1$, $n=1$ and $e_L=e_F$. \Cref{prop:general} is the relevant benchmark for a system with those terms; whether the deployed system satisfies its other hypotheses is not something we have established, and we do not claim that it is an instance of the model. In every borrow event in our data the fee is $\EFeeMedian\%$ of the gross amount. Two composite operations exist: \emph{leverage} borrows and buys tokens in one transaction, and \emph{deleverage} sells collateral and repays in one transaction, which is the collateral-netted close of \cref{prop:netclose} applied to all or part of a position.

We collected every event log of the system's contract on Ethereum from its deployment (block \EStartBlock) to block \EEndBlock: \ELogs{} logs over \EDays{} days, covering \EPools{} pools. Events were decoded with the event definitions of the project's public software kit (version 1.3.0); every log is either decoded, or belongs to a listed set of administrative and fee-distribution event types, or is the one log of an event type that we could not identify, whose data is a single address and which we exclude. A separate ledger decodes the fee-distribution events and reward claims from the supplied ABI definitions (\cref{sec:rewards}). An \emph{account} is a pair of a pool and an address. The collateral and debt of each account are reconstructed from the events that carried accounts over from the previous version at launch and from the post-state recorded in every later credit event; staked balances are reconstructed from the staking events, and the analysis checks that in every pool the reconstructed stake covers the collateral and does not exceed the contract's token balance. Each stored log carries the emitting address, the block hash, the transaction hash, the transaction index and the log index, so that the sample can be checked against any archive node. At the last block we also read the contract's own state for each pool: floor price, pool price, total collateral, total debt, total stake and circulating supply. The reconstructed collateral, debt and stake of every pool equal the contract's totals exactly, and the analysis stops if they do not. All six pools use the same reserve asset (wrapped ether). Durations in days are computed from block timestamps, interpolated linearly between the timestamps of fetched anchor blocks. Time windows are measured in blocks: \EHourBlocks{} blocks are about one hour and \EWeekBlocks{} blocks about seven days. The scripts and the raw logs are in the supplementary material (\cref{app:repro}).

Four limits apply to everything below. Most accounts were carried over from the previous version and show no later activity, so inactivity and deliberate deferral cannot be told apart for them. The floor did not change in \EFlatPools{} of the \EPools{} pools during the window. We observe addresses, not people, and an event that names an account need not have been initiated by its owner. And the pools belong to one protocol and are not independent observations.

\subsection{Results}

\begin{table}[t]
\centering\small
\caption{Baseline credit on Ethereum by pool, block \EStartBlock{} to \EEndBlock. All state variables are read at the last block. ``Locked share'' is collateral divided by the circulating supply that the contract reports, which excludes the pool's own inventory. ``Debt/floor'' is total debt divided by the value of total collateral at the floor price. ``Premium'' is the pool price divided by the floor at the first and last swap. ``Floor changes'' counts changes of the floor between consecutive swaps. ``Incr.'' counts borrow and leverage events, including increases of existing loans. Command: \texttt{python3 -I empirics/analyze.py}.}\label{tab:baseline}
\setlength{\tabcolsep}{4pt}
\begin{tabular}{@{}lrrrrrrrrrrrr@{}}
\toprule
 & \multicolumn{2}{c}{Accounts} & Locked & Debt/ & \multicolumn{2}{c}{Premium} & Floor & \multicolumn{2}{c}{Carried over} & \multicolumn{3}{c}{Credit events}\\
\cmidrule(lr){2-3}\cmidrule(lr){6-7}\cmidrule(lr){9-10}\cmidrule(lr){11-13}
Pool & ever & in debt & share & floor & first & last & changes & all & in debt & incr. & repay & delev.\\
\midrule
\EPoolRows
\bottomrule
\end{tabular}
\end{table}

\paragraph{The loan book is large relative to circulating supply.} In pool B, which has the most accounts, \EBShare\% of the circulating supply is locked as collateral, in \EBOpen{} accounts, and total debt is \EBLtv\% of the floor value of the collateral (\cref{tab:baseline}). Across the six pools the locked share ranges from \EShareMin\% to \EShareMax\%. The denominator matters in one pool. For ONYX the contract reports a total supply of \EOnyxLensTotal{} tokens, while the token's total supply on Ethereum is \EOnyxErcTotal; the difference of \EOnyxOff{} tokens is part of the circulating supply of \EOnyxCirc{} that the contract reports. We conjecture, and have not confirmed, that these are tokens on other chains. If they are left out of the denominator, the locked share of ONYX is \EOnyxShareErc\%. For comparison, \cref{thm:eq} gives a locked share of $98\%$ for an origination fee of $1\%$, an outside rate of $5\%$ and exponentially distributed horizons with a mean of one year. The horizon distribution and the outside rates are not observed, and most positions were inherited from the previous version, so this is an illustration and not a test. Most of the circulating supply of most pools is borrowed against. The separate debt/floor ratio conditions on the collateral already locked: if a facility sizes the lock to support the requested debt, this ratio can be close to one mechanically. It does not measure whether an account borrowed against all its eligible holdings. That requires a history of eligible balances, including unlocked tokens, which we do not reconstruct.

\paragraph{Plain repayments are mostly followed by sales.} There are \EReduce{} debt-reducing events. Of these, \EDeleverage{} (\EDelevPct\%) are deleverage operations, in which a sale of collateral and the repayment are one transaction. They coincide with a sale by construction and are evidence about how loans are closed, not about why: of the full closures, \EFullDelev{} were made by a deleverage and \EFullRepay{} by a plain repayment, so the netted close of \cref{prop:netclose} is the usual way out.

The informative events are the \ERepay{} plain repayments, after which the borrower is free to keep the released tokens. \ECatSold{} of them (\ERepaySoldPct\%) are followed within \EWeekBlocks{} blocks by a sale of tokens by the same account, \ESoldHour{} of them within \EHourBlocks{} blocks; \ECatReborrow{} are followed by a new loan and no sale, and \ECatNeither{} by neither. Every repayment has a complete follow-up window inside the sample. The events are concentrated: two accounts in the ONYX pool made \ETopTwoEvents{} of the repayments (\ETopTwoPct\%), and without them the share followed by a sale is \ERestPct\% (\ERestSold{} of \ERestN). Counting accounts, \EAcctSoldAfter{} of the \ERepayAccts{} accounts that made a plain repayment sold within \EWeekBlocks{} blocks after one. Weighting by the amount repaid, \EAmtSoldPct\% is followed by a sale. The share is the same for full repayments (\EFullRepSold{} of \EFullRepN) and partial ones (\EPartRepSold{} of \EPartRepN), and it is \ENoOnyxSold{} of \ENoOnyxN{} without the ONYX pool. For a base rate, take the same accounts at an arbitrary date in the sample: weighting each account by its number of repayments, a sale by the account falls in \EBaseRate\% of week-long windows. A sale is thus more than twenty times as frequent after a repayment as at a random date. The comparison shows that repayments and sales cluster in time. It does not show that the repayment is made to sell, since both can be parts of one decision to leave. Resampling the \ERepayClusters{} accounts, not the events, gives a 95\% interval of \ERepayCILo\% to \ERepayCIHi\% for the share of events (\EAddrCILo\% to \EAddrCIHi\% when the \EAddrClusters{} addresses are resampled, an address being able to hold accounts in several pools); the interval describes the dependence on which accounts are in the sample and is not a test of a model. The estimate is not uniform (\cref{tab:repaypool}). It is \ERegOneSold{} of \ERegOneN{} up to the last upgrade of the credit component and \ERegTwoSold{} of \ERegTwoN{} after it, \EHalfOneSold{} of \EHalfOneN{} in the first half of the window and \EHalfTwoSold{} of \EHalfTwoN{} in the second, and in pool B it is one half. We do not know whether the change over time reflects the upgrade, the fall of the premium, or the two ONYX accounts, whose repayments are late in the window. We cannot tell whether the tokens sold are the tokens released. A window before the repayment is not a clean placebo, because accounts unwind in repeated steps of repaying and selling: \ESoldBeforeHour{} repayments are also preceded by a sale within \EHourBlocks{} blocks. As a rough comparison over a much longer window, of the \ENeverReduced{} accounts that were in debt at the end and never reduced their debt, \ENeverReducedSold{} (\ENeverReducedSoldPct\%) made any sale at all. Between half and four fifths of plain repayments, depending on the weighting, are thus followed by a sale. That is the direction \cref{prop:repay} predicts, on a small and concentrated sample, and it does not establish the mechanism.

\begin{table}[t]
\centering\small
\caption{Plain repayments by pool, and how many are followed within \EWeekBlocks{} blocks by a sale by the same account.}\label{tab:repaypool}
\begin{tabular}{@{}lrrr@{}}
\toprule
Pool & Repayments & Accounts & Followed by a sale\\
\midrule
\ERepayPoolRows
\bottomrule
\end{tabular}
\end{table}

\paragraph{Loans stay open.} Of the \ESnap{} accounts carried over at launch, \ESnapOpen{} (\ESnapOpenPct\%) were still in debt at the end of the window and \ESnapSame{} (\ESnapSamePct\%) had the same collateral and debt as at launch. \EActive{} of the carried-over accounts appear in some later event (a credit operation, a swap, a reward claim, a withdrawal of stake, or a deposit of stake more than ten blocks after the carry-over, which excludes the deposits that belong to the migration itself), and of those \EActiveClosed{} (\EActiveClosedPct\%) closed their loan. Most open loans are thus loans of accounts that show little or no activity. That is consistent with \cref{prop:repay}, under which a borrower who does not want to sell has no reason to act; it is equally consistent with inattention or lost access, and these data cannot separate the explanations.

\paragraph{New accounts mostly enter through leverage, and many episodes are short.} Gross new debt in the window is \EAmtLever{} units of the reserve asset from leverage operations and \EAmtBorrow{} from plain loans, a leverage share of \ELeverAmtPct\%. \ENewAccts{} accounts that were not carried over opened \EEpisodes{} loan episodes in the window; \EEpLeverage{} episodes were opened by a leverage operation and \EEpBorrow{} by a plain loan. \EEpClosed{} episodes were closed in the window, at a median age of \EEpMedianClose{} days, \EEpClosedThirty{} of them within thirty days; these short episodes belong to \EAcctsShort{} of the accounts. The other \EEpOpen{} episodes were still open at the end of the window, at a median age of \EEpMedianOpenAge{} days. Episodes that begin late in the window are censored, so the closed episodes overstate how short loans are. A Kaplan--Meier estimate, which counts each episode once and assumes that the date an episode begins is unrelated to its length, puts the share of new episodes still open at \EKMThirty\% after 30 days and at \EKMNinety\% after 90 and after 180 days. About half of new loan episodes end within a month. The numbers of episodes still under observation at 30, 90 and 180 days are \EKMRiskThirty, \EKMRiskNinety{} and \EKMRiskHalfYear. The estimated curve is flat between days 90 and 180, and \EKMLateCloses{} episode in the sample closes later than that. For the first episode of each new account alone the estimates are \EKMFirstThirty\%, \EKMFirstNinety\% and \EKMFirstHalfYear\%. An episode is the life of a loan and not the time for which its owner holds the token, so these durations do not estimate the horizons $T$ of \cref{sec:equilibrium}. Of the first episodes of the \ENewAccts{} accounts, \EFirstClosed{} were closed and \EFirstClosedThirty{} within thirty days. A loan held for a few days at a fee of $1\%$ is far inside the break-even horizon of $2.4$ months that \cref{prop:util} gives for a holder who borrows for cash at an outside rate of $5\%$. Borrowers who enter through leverage borrow to buy more of the token. That motive is outside the assumptions of \cref{sec:incentives,sec:equilibrium}, so their behavior does not bear on those results either way. It is a use of credit outside the benchmark. We do not observe beliefs, outside rates or account-specific accrued rewards, so the data do not test \cref{cor:noentry} and do not identify what brings buyers in.

\paragraph{The floor rose in one pool only.} We observe the floor at every swap and in the contract's state at the last block, not continuously. In \EFlatPools{} of the \EPools{} pools the floor was the same at the last swap as at the first and at every swap between, and in those pools the floor read from the contract at the last block equals the floor at the last swap, which lies up to \EMaxSwapGap{} days before the end. In the same pools the premium of the pool price over the floor shrank; in pool B it went from \EBPremFirst{} to \EBPremLast. A floor that cannot fall is not a floor that rises: as \cref{sec:raise} makes explicit, it rises only when reserve is paid in and applied to it. In the sixth pool, ELEFENT, the floor changed \EElChanges{} times and rose by a factor of \EElRatio. That pool has only \EElAccounts{} credit accounts, \EElOpen{} of them in debt, and their debt is \EElLtv\% of the value of their collateral at the floor of the last block, which is \EElFloorEndPct\% above the floor at the last swap. The ratio is lower than in any other pool: these few borrowers had not borrowed all the capacity that the rising floor gave them. Taken at face value this is evidence against the top-up plan of \cref{thm:eq}(a), under which a borrower with a long horizon draws every unit of new capacity. It concerns five accounts in one pool, and we do not know their horizons or whether the fee makes small top-ups unattractive; the model has no fixed cost of a transaction. The sample is too small to say more about top-ups, and we did not examine minting ahead of the floor changes in this pool.

\subsection{Fee allocations and reward claims}\label{sec:rewards}

The same logs contain \RNDistributions{} \texttt{Distributed} events, whose ABI records a total amount and its allocation to the protocol, creator and staking. We decode these fields directly from the supplied ABI and require, for every event, that the three allocations sum exactly to the total. The ledger records \RTotal{} WETH allocated, including \RStaking{} WETH to staking (\cref{tab:rewards}). These are fee destinations recorded by events, not measured earnings of particular accounts. The table reports the separate \RNClaims{} \texttt{Claim} events, totaling \RClaimed{} WETH, without assigning them to the distributions of this window. Opening accruals, each account's reward history and the reward-accounting rules have not been reconstructed. Subtracting claims from allocations therefore does not identify an unclaimed balance, and neither column is an estimate of a holder's yield.

\begin{table}[t]
\centering\small
\caption{Recorded fee allocations and reward claims, in WETH, in the sample of \cref{tab:baseline}. The four allocation columns reconcile exactly in integer units; displayed amounts are rounded. Claims are a separate series and need not settle rewards allocated in the same window. Command: \texttt{python3 -I empirics/rewards.py}.}\label{tab:rewards}
\begin{tabular}{@{}lrrrrr@{}}
\toprule
Pool & Total allocated & Protocol & Creator & Staking & Claimed\\
\midrule
\RPoolRows
\bottomrule
\end{tabular}
\end{table}

An allocation is assigned to a fee source only when the next borrow, leverage or swap event for the same pool in the same transaction supplies an exactly matching amount, with each source used once. The matching amounts are the recorded borrow fee, the debt added less reserves used in a leverage, and the swap fee less its liquidity allocation. All \RCreditN{} origination events match \RCreditTotal{} WETH of distributions, of which \RCreditStaking{} WETH is allocated to staking; all \RSwapN{} swap events match \RSwapTotal{} WETH. The remaining \RUnassignedN{} distributions, totaling \RUnassignedTotal{} WETH, retain their recorded allocations but are left unassigned. These matches reconcile event accounting; they do not establish who ultimately bears or receives the fees. In particular, reserves used in a leverage are borrowed funds spent in its composite purchase, not a separate cash contribution by the user.

This ledger measures a payout channel excluded by (A2) of the basic incentive benchmark and admitted by the extension in \cref{sec:general}. Since borrowers are eligible for staking rewards, part of the fees generated by credit may flow to borrowers as members of that group. The data do not identify their share, their net financing costs, or whether rewards cause borrowing. Establishing those quantities requires account reward accruals, opening inventories and the complete asset transfers of each transaction. The allocations also differ across pools, but six pools in one protocol do not provide a causal comparison of reward policies.

\subsection{Assessment}

Two observations are consistent with the theory under its assumptions. At a low origination fee most of the circulating supply of most pools is borrowed against. And most plain repayments are followed by a sale by the same account, although two accounts supply almost half of those events, the share is one half in the pool with the most accounts, and it differs between the early and the late part of the window. The prevalence of the netted close shows which operation borrowers use and is not evidence about motive. One observation shows that the model does not describe this market: nine tenths of new debt finances purchases of the token, and about half of new loan episodes end within a month. The borrower of \cref{sec:incentives}, who wants cash and holds for longer than the break-even horizon, is the minority among new borrowers here. The fee ledger further shows an aggregate reward channel that the no-payout benchmark omits, without identifying individual earnings. The statements about extraction and exit routing were not examined; the sell fee at the floor is close to zero by design, and the floor moved in one small pool only. A study across protocols, over a period in which floors move, with a model of leveraged demand and fee redistribution to guide it, is the natural next step.
